% Bewertungspaket zum Gesamttest «Lineare Funktionen» — Quelle des PDFs (python3 scripts/build-lp-pdf.py).
% Ganze Punkte je Teilschritt; Werte mit python3 nachgerechnet (03.10.2026, Folgefehler-Fälle eingeschlossen).
\documentclass[11pt]{article}
\usepackage{../lp-druck}
\renewcommand{\lpname}{Lineare Funktionen}
\renewcommand{\lpteilgebiet}{3.2}
\renewcommand{\lpfuss}{Bewertungspaket}
\newcolumntype{P}{>{\centering\arraybackslash}p{10mm}}
\newenvironment{raster}{\par\smallskip\noindent\tabularx{\linewidth}{@{}>{\raggedright\arraybackslash}p{38mm}P X@{}}\toprule
  \small\textsc{Teilschritt} & \small\textsc{P} & \small\textsc{Musterlösung}\\\midrule}{\bottomrule\endtabularx\par}
\newcommand{\fehler}[1]{\par\smallskip{\small\textcolor{lprot}{\bfseries Typische Fehler:} #1}}
\newcommand{\E}{\,{\footnotesize\bfseries\color{lpgruen}(E)}}
\newcommand{\A}{\,{\footnotesize\bfseries\color{lpblau}(A)}}
\begin{document}
\lpkopf{Bewertungspaket zum Gesamttest}{}

\begin{lpachtung}{Erst lösen, dann öffnen}
Dieses Paket enthält die Musterlösung. Wer es vor dem Lösen liest, misst am Ende nur, wie gut das Abschreiben gelingt.
\end{lpachtung}

\section*{1 · So gehst du vor}
\begin{enumerate}[leftmargin=*,itemsep=2pt]
\item \textbf{Gesamttest lösen} — vollständig, mit Rechenweg, ohne dieses Paket und ohne Taschenrechner.
\item \textbf{Fotografieren oder scannen}, gut lesbar, eine Seite pro Bild. \textbf{Kein Name} auf den Bildern.
  Handyfotos tragen oft unsichtbar Standort, Zeit und Gerät mit (Metadaten). Lade darum lieber einen Scan oder ein
  Bildschirmfoto des Fotos hoch, oder entferne vorher die Standortangaben.
\item \textbf{Bewerten lassen.} Ob du dafür eine KI nutzt und welche, entscheidest du selbst und in eigener Verantwortung —
  je nachdem, was dir aufgrund deines Alters und deiner persönlichen Situation erlaubt ist (Altersgrenzen und
  Nutzungsbedingungen des Anbieters, allenfalls Einverständnis deiner Eltern). Lade nur deine Lösung und dieses PDF hoch,
  keine persönlichen Daten, und füge den Auftrag aus Abschnitt~2 ein. Ohne KI kannst du dich mit Abschnitt~3 selbst bewerten.
\item \textbf{Rückmeldung prüfen.} Stimmt, was die KI aus deiner Handschrift gelesen hat? Eine KI kann sich irren —
  im Zweifel gilt das Raster in Abschnitt~3.
\item \textbf{Gezielt wiederholen} nach Abschnitt~4.
\end{enumerate}

\needspace{14\baselineskip}
\section*{2 · Auftrag an die KI}
\begin{tcolorbox}[colback=lplinie!25,colframe=lplinie,boxrule=0.5pt,arc=1mm,fontupper=\small]
Du bist Korrektorin oder Korrektor für Mathematik an einer Schweizer Berufsmaturitätsschule. Im Anhang findest du (1)~das Bewertungspaket zum Gesamttest «Lineare Funktionen» mit Musterlösung und Punkteraster und (2)~Fotos meiner handschriftlichen Lösung.\par\medskip
Geh so vor:\par
1. Schreib zu jeder Aufgabe G1 bis G8 zuerst ab, was du in meiner Lösung liest — Rechenweg und Ergebnis. Was du nicht sicher lesen kannst, markierst du mit [unsicher]. Nicht raten.\par
2. Vergib die Punkte genau nach dem Raster im Bewertungspaket (Abschnitt 3), ganze Punkte je Zeile. Ziehe einen Fehler nur einmal ab: Rechne ich mit einem falschen Zwischenergebnis richtig weiter, gibt es die Folgepunkte. Die Zeilen «Typische Fehler» sagen, welche Rasterzeile bei diesem Fehler 0 Punkte gibt — sie sind kein zusätzlicher Abzug.\par
3. Andere richtige Lösungswege zählen voll. Gleichwertige Schreibweisen zählen voll: Dezimalpunkt oder Dezimalkomma, Bruch oder Dezimalzahl ($\tfrac23 \approx 0.67$, $-\tfrac12 = -0.5$), $y = \ldots$ oder $f(x) = \ldots$, umgestellte Terme ($3 - 1.5x$ statt $-1.5x + 3$) — ausser wo das Raster eine bestimmte Form verlangt. Zeilen mit (E) sind Ergebnispunkte, Zeilen mit (A) verlangen nur Ablesen: Diese Punkte gibt es auch ohne Weg. Alle anderen Punkte gibt es nur mit nachvollziehbarem Weg.\par
4. Begründe jeden Abzug in einem Satz und nenne die Fehlerart (zum Beispiel «Vorzeichen bei $b = y_1 - m\,x_1$»). Schreib die Musterlösung nicht ab — zeig mir, wo mein Fehler liegt.\par
5. Gib zum Schluss eine Tabelle «Aufgabe | Punkte | Begründung», die Gesamtpunktzahl (von 22) und — nach der Selbsteinschätzung in Abschnitt 4 — welches Kapitel des Leitprogramms ich wiederholen soll. Keine Note.\par
6. Fehlt ein Foto oder ist eine Aufgabe unlesbar, sag es, statt sie zu bewerten.
\end{tcolorbox}

\needspace{16\baselineskip}
\section*{3 · Musterlösung und Punkteraster}
\textbf{Allgemein:} Ganze Punkte je Zeile. Ein Fehler wird nur einmal abgezogen (Folgefehler). Gleichwertige Wege und
Schreibweisen zählen voll (Bruch oder Dezimalzahl, Terme in anderer Reihenfolge). In der Spalte P: \textbf{(E)}~= Ergebnispunkt,
\textbf{(A)}~= nur ablesen: Diese Punkte gibt es auch ohne Weg, alle anderen nur mit nachvollziehbarem Weg.
Der ganze Test ist ohne Taschenrechner gestellt. «Typische Fehler» nennt, welche Zeile 0 Punkte gibt, und die Punkte, die bleiben.

\begin{lpaufgabe}{G1}{$f(x) = -1.5x + 3$ zeichnen}{3 P}
\begin{raster}
Zwei Punkte bestimmt & 1 & z.\,B. $(0 \mid 3)$ und $(2 \mid 0)$ — oder $b = 3$ und von dort 2 nach rechts, 3 hinunter\\
Gerade gezeichnet & 1\E & Gerade durch beide Punkte, über das ganze Fenster gezogen\\
Nullstelle & 1\E & $0 = -1.5x + 3 \Rightarrow x_0 = 2$\\
\end{raster}
\fehler{Gerade steigt statt fällt (Vorzeichen von $m$): Gerade 0; die Nullstelle $x_0 = 2$ zählt trotzdem $\to$ 2 von 3\,P.
Nur die beiden Punkte markiert, keine Gerade gezogen: Gerade 0 $\to$ 2 von 3\,P.
$x_0 = -2$ (Vorzeichen beim Auflösen): Nullstelle 0 $\to$ 2 von 3\,P.
Nullstelle nur abgelesen, ohne Rechnung: zählt voll — sie ist im Bild sichtbar.}
\end{lpaufgabe}

\begin{lpaufgabe}{G2}{Graph $\to$ Gleichung}{3 P}
\begin{raster}
$b$ abgelesen & 1\A & $b = -2$, denn die Gerade schneidet die $y$-Achse bei $-2$\\
$m$ aus dem Steigungsdreieck & 1 & von $(0 \mid -2)$ nach $(3 \mid 0)$: 3 nach rechts, 2 hinauf, also $m = \tfrac23$\\
Gleichung und $f(-3)$ & 1\E & $f(x) = \tfrac23 x - 2$; $f(-3) = \tfrac23 \cdot (-3) - 2 = -4$\\
\end{raster}
\fehler{$m = \tfrac32$ (Bruch verkehrt): $m$-Zeile 0; $f(-3) = -6.5$ folgerichtig $\to$ 2 von 3\,P.
$m = -\tfrac23$ (Richtung verwechselt): $m$-Zeile 0; $f(-3) = 0$ folgerichtig $\to$ 2 von 3\,P.
$b = 3$ (Nullstelle statt Achsenabschnitt): $b$-Zeile 0, Rest folgerichtig $\to$ 2 von 3\,P.
$f(-3)$ fehlt oder ist falsch gerechnet, obwohl die Gleichung stimmt: letzte Zeile 0 $\to$ 2 von 3\,P.
$-3$ einfach am Graphen abgelesen ($-4$): zählt voll, der Punkt liegt auf einem Gitterpunkt.}
\end{lpaufgabe}

\begin{lpaufgabe}{G3}{$f(x) = 0.5x - 3$}{3 P}
\begin{raster}
(a) $b$ und Nullstelle & 1 & $b = -3$; $0 = 0.5x - 3 \Rightarrow x_0 = 6$\\
(b) Punktprobe & 1\E & $f(10) = 0.5 \cdot 10 - 3 = 2$: $P(10 \mid 2)$ liegt auf dem Graphen\\
(c) Deutung von $m$ & 1 & Die Gerade steigt: pro Schritt nach rechts geht es $0.5$ hinauf (oder: pro 2 nach rechts 1 hinauf).\\
\end{raster}
\fehler{$x_0 = -6$ (Vorzeichen) oder $x_0 = -3$ ($b$ als Nullstelle): Zeile (a) 0 $\to$ 2 von 3\,P.
Nur «ja» ohne Rechnung bei (b): Zeile (b) 0 — hier ist der Weg verlangt $\to$ 2 von 3\,P.
(c) «die Gerade ist flach» ohne Aussage über den Zuwachs: Zeile (c) 0. «steigt, weil $m > 0$» zählt voll.}
\end{lpaufgabe}

\begin{lpaufgabe}{G4}{$g: y = -4x + 1$}{3 P}
\begin{raster}
(a) Parallele & 1\E & gleiche Steigung, $b = -3$ ist direkt gegeben: $y = -4x - 3$\\
(b) Steigung senkrecht & 1\E & $m_1 \cdot m_2 = -1 \Rightarrow m_2 = -\tfrac{1}{-4} = 0.25$\\
(c) Begründung & 1 & $x = 5$ ordnet der einen Stelle $x = 5$ unendlich viele $y$-Werte zu; eine Funktion ordnet jedem $x$ genau einen Wert zu.\\
\end{raster}
\fehler{$y = -4x + 1$ bei (a) (dieselbe Gerade statt einer parallelen durch den gegebenen Punkt): Zeile (a) 0 $\to$ 2 von 3\,P.
$m_2 = 4$ oder $m_2 = -0.25$ (Kehrwert oder Vorzeichen fehlt): Zeile (b) 0 $\to$ 2 von 3\,P.
(c) «weil sie senkrecht steht» ohne den Grund (ein $x$, viele $y$): Zeile (c) 0. Der Hinweis «$m$ ist nicht definiert, weil $\Delta x = 0$» zählt voll.}
\end{lpaufgabe}

\begin{lpaufgabe}{G5}{Drei Fälle}{2 P}
\begin{raster}
Zuordnung & 1\A & $y = 0.5x$ proportional ($b = 0$) · $y = -2$ konstant ($m = 0$) · $x = 4$ keine Funktion\\
Begründung & 1 & $x = 4$ ist eine senkrechte Gerade: zu dieser einen Stelle gehören unendlich viele $y$-Werte.\\
\end{raster}
\fehler{$y = -2$ und $x = 4$ vertauscht: Zuordnung 0; eine folgerichtige Begründung zu $y = -2$ gibt den zweiten Punkt nicht $\to$ 0 von 2\,P.
Zuordnung richtig, Begründung fehlt: 1 von 2\,P.}
\end{lpaufgabe}

\begin{lpaufgabe}{G6}{$A(-3 \mid 6)$, $B(2 \mid -4)$}{3 P}
\begin{raster}
Steigung & 1 & $m = \dfrac{-4 - 6}{2 - (-3)} = \dfrac{-10}{5} = -2$\\
$b$ und Gleichung & 1\E & $6 = -2 \cdot (-3) + b \Rightarrow b = 0$, also $f(x) = -2x$\\
Typ & 1 & $b = 0$: eine proportionale Funktion — die Gerade geht durch den Ursprung.\\
\end{raster}
\fehler{$m = 2$ (Reihenfolge im Bruch vertauscht): Steigung 0; $b = 12$ und $f(x) = 2x + 12$ folgerichtig, Typ dann nicht proportional $\to$ 1 von 3\,P.
$m = -10$ ($\Delta y$ ohne Teilen): Steigung 0, Rest folgerichtig $\to$ 2 von 3\,P (Typ nur, wenn das gefundene $b$ wirklich 0 ist).
$b = 12$ ($m\,x_1$ dazugezählt statt abgezogen, $6 + (-2) \cdot (-3)$): $b$-Zeile 0, Typ 0 $\to$ 1 von 3\,P.
Gleichung richtig, Typ fehlt oder «linear»: Typ 0 $\to$ 2 von 3\,P.}
\end{lpaufgabe}

\begin{lpaufgabe}{G7}{$m = -0.5$, $P(6 \mid 1)$}{2 P}
\begin{raster}
$b$ bestimmt & 1 & $1 = -0.5 \cdot 6 + b \Rightarrow b = 1 + 3 = 4$\\
Gleichung & 1\E & $f(x) = -0.5x + 4$ (gleichwertig $-\tfrac12 x + 4$)\\
\end{raster}
\fehler{$b = -2$ ($m\,x_1$ dazugezählt statt abgezogen): $b$-Zeile 0; $f(x) = -0.5x - 2$ folgerichtig $\to$ 1 von 2\,P.
Nur $b = 4$ ohne die Gleichung: Gleichung 0 $\to$ 1 von 2\,P.}
\end{lpaufgabe}

\begin{lpaufgabe}{G8}{Tank, 300 Liter, 12 Liter je Minute}{3 P}
\begin{raster}
Funktionsgleichung & 1 & «pro Minute» $\to m = -12$ (es fliesst ab), «enthält 300 Liter» $\to b = 300$: $V(t) = -12t + 300$\\
Leer nach & 1\E & $0 = -12t + 300 \Rightarrow t = 25$ Minuten\\
Definitionsmenge & 1 & $0 \le t \le 25$ — vorher gibt es keine Messung, nachher ist der Tank leer (gleichwertig $D = [0;\,25]$).\\
\end{raster}
\fehler{$V(t) = 12t + 300$ (Vorzeichen): Gleichung 0; es gibt keine Nullstelle, «leer nach» 0, eine folgerichtige Definitionsmenge $t \ge 0$ gibt den dritten Punkt $\to$ 1 von 3\,P.
$V(t) = 300t - 12$ (Zahlen vertauscht): Gleichung 0, $t = 0.04$ folgerichtig $\to$ 1 von 3\,P (Definitionsmenge nur, wenn sie zum eigenen Ergebnis passt).
$D = \mathbb{R}$ oder keine Angabe: Definitionsmenge 0 $\to$ 2 von 3\,P. $0 < t < 25$ oder $t \ge 0$ mit Hinweis auf die Leerung zählt voll.}
\end{lpaufgabe}

\needspace{14\baselineskip}
\section*{4 · Selbsteinschätzung}
\begin{tabularx}{\linewidth}{@{}l X@{}}\toprule
19 – 22 P & Die geprüften Teile sitzen. Wo du Punkte verloren hast: das Kapitel dieser Aufgabe nochmals (Zuordnung unten).\\
15 – 18 P & Den schwächsten Teil nochmals: Simulation und Übungen des Kapitels, in dem du die meisten Punkte verloren hast.\\
10 – 14 P & Zurück zu den Kapiteln aller Aufgaben, in denen du Punkte verloren hast.\\
0 – 9 P & Zurück zu Kapitel 1 und von dort der Reihe nach weiter.\\\midrule
\multicolumn{2}{@{}p{\linewidth}@{}}{\small\textbf{Aufgabe $\to$ Kapitel:} G1 $\to$ Kapitel 1 (zeichnen) und 2 (Nullstelle) · G2 $\to$ Kapitel 1 und 2 · G3 $\to$ Kapitel 1 und 2 · G4 $\to$ Kapitel 3 (Lage) und 4 (aufstellen) · G5 $\to$ Kapitel 3 · G6, G7, G8 $\to$ Kapitel 4}\\\bottomrule
\end{tabularx}
\end{document}
